\section{通信效率优化方法}\label{sec:methods}

由式~\eqref{eq:total-bits}和式~\eqref{eq:wireless-link}可知，通信开销可以从参与节点、交互频率、单次载荷、无线传输过程和消息路径五个层面处理。下文以方法直接改变的变量为分类依据，而不按算法名称逐篇介绍。

\subsection{客户端选择与通信频率控制}\label{subsec:selection-frequency}

移动边缘训练中，可用设备常多于基站能够同时接入的数量，部分设备还会离线或逾期。客户端选择决定参与集合 \(\mathcal S_t\)，频率控制决定本地步数 \(E_k\)、触发条件和更新时序。两者分别回答“选谁”和“何时通信”，影响参与规模、全局轮数或等待时间。

客户端选择分为随机、资源感知和数据效用感知三类。随机选择无需详细状态，但可能选中慢设备。资源感知选择根据算力、电量、在线状态和信道质量，筛选能按期完成训练与上传的设备。数据效用感知还估计样本覆盖、局部损失或更新价值\cite{Nishio2019FedCS,Yang2020Scheduling,Xu2021ClientBandwidth,Zhang2024ClientSelectionSurvey}。“快设备优先”虽能缩短等待，却会产生选择偏差：设备性能和位置可能与数据分布相关，长期排除弱设备会使参与数据偏离总体。

通信时序可分为同步、异步和事件触发机制。同步训练等待本轮更新到齐后再聚合，轮次时间受最慢参与者限制。异步训练允许更新到达后立即或分批采用，主要减少同步空等。更新陈旧度是客户端所用模型版本与服务器当前版本的差距，差距越大，更新越可能偏离当前方向。异步不必然减少累计上传；若每次更新仍单独发送，它减少的是等待而非比特。事件触发只在模型变化或梯度价值超过阈值时上传，才可能过滤冗余通信\cite{Lu2020EAFLM,Liu2021AsyncFL}。

相关研究已从“让本轮按期完成”扩展到长期公平和动态到达。FedCS按资源状态选择设备，后续无线调度与长期选择再考虑信道、带宽和跨轮参与；自适应方法则按资源预算调整聚合周期\cite{Nishio2019FedCS,Yang2020Scheduling,Xu2021ClientBandwidth,McMahan2017FedAvg,Wang2019AdaptiveFL}。异步边缘学习通过阈值或部分聚合抑制低价值更新，通信—计算重叠与公平调度用于减少等待并保留弱设备机会\cite{Lu2020EAFLM,Liu2021AsyncFL,Ozfatura2021FedTOE}。这些机制分别处理逾期参与、同步空等或冗余上传，并不同时降低所有指标。

这类方法需要状态上报和调度计算，还可能产生陈旧更新与公平性问题。随机选择适合设备条件相近的系统；资源感知选择适合期限严格的场景；数据效用感知适合数据差异明显且效用可估计的任务。异步机制适合设备间歇在线，但不适合对陈旧更新敏感的训练。参与集合确定后，设备仍要共享频谱，下一步还需配置无线与计算资源。

\subsection{无线资源分配与联合优化}\label{subsec:resource-allocation}

摄像头、车辆或手机共同训练时，弱信道或低算力终端仍会拖慢整轮训练。同步聚合需等待更新到齐，最慢终端的“本地计算时间加上传时间”常决定轮次完成时间。资源分配配置带宽 \(b_{k,t}\)、功率 \(P_{k,t}\)、接入时隙、中央处理器（CPU）频率和边缘关联，以改善传输时延、能耗或掉队率。

这些变量解决的问题不同。带宽决定频谱份额，功率影响接收质量与速率，时隙决定发送时间；三者作用于无线传输。CPU频率改变本地计算速度和计算能耗，不直接进入上行速率式。边缘关联决定终端连接哪个基站或边缘服务器，进而影响链路质量、服务器负载和辅助计算位置。系统通常先收集信道、能量和计算状态，再在时延或能量约束下求解配置\cite{Tran2019WirelessOptimization,Dinh2021WirelessConvergence}。这些操作通常不减少 \(B_{k,t}^{\mathrm{ul}}\)，主要改变传完相同比特所需的时间和能量。

资源决策分为单轮配置和跨轮配置。单轮方法依据当前状态控制本轮时延、能耗或掉队数量，适合状态变化较慢且要求快速求解的场景。跨轮方法还考虑剩余电量、历史参与和未来信道，避免持续透支同一批设备。相关研究由此扩展到收敛约束、信道调度和长期客户端—带宽设计\cite{Tran2019WirelessOptimization,Dinh2021WirelessConvergence,Yang2020Scheduling,Xu2021ClientBandwidth}。

多个联邦任务共享接入网时，资源配置还要在业务之间分配频谱。边缘辅助方案会同时决定弱终端是否由边缘侧辅助，并配置CPU频率、无线资源和模型稀疏度\cite{Xu2022MultiFLBandwidth,Chen2025EdgeAssistedSparsification}。这类联合决策能同时处理通信等待、计算等待和长期能量约束，但变量增多会提高在线求解难度。在线学习或强化学习虽能适应动态状态，也会产生探索、训练和控制开销。

资源优化适合蜂窝边缘、多业务接入和设备差异明显的系统。若网络稳定、带宽充足或状态难以获取，复杂优化的收益可能不足以抵消开销。信道反馈占用控制比特，在线求解消耗时间和算力，求解期间状态变化还可能使配置过期。因此，反馈和调度应计入 \(B_{k,t}^{\mathrm{ctl}}\)，比特、时延和能耗也应分别报告。资源分配解决“怎样把更新传完”；若更新本身过大，还需用压缩削减载荷。

\subsection{模型压缩与通信载荷削减}\label{subsec:model-compression}

深度模型参数很多，即使信道较好，完整上传模型或梯度仍会占用大量带宽。模型压缩直接改变单轮载荷：量化减少每个数值使用的比特数，稀疏化减少需要发送的元素数量。量化先把连续数值映射到有限等级，再发送等级编码；稀疏化先按幅值或阈值保留少量重要元素，其余元素本轮不发送。两者可以组合，但会分别引入量化误差和丢弃误差。

为区分数学稀疏性与真实传输开销，设待传更新为 \(\bm g_{k,t}\in\mathbb R^d\)，稀疏算子为 \(S(\cdot)\)，本文定义
\begin{equation}\label{eq:effective-compression}
  \begin{aligned}
    \sigma_{k,t}
      &=1-\frac{\lVert S(\bm g_{k,t})\rVert_0}{d},\\
    \rho_{k,t}^{\mathrm{eff}}
      &=\frac{B_{k,t}^{\mathrm{val}}+B_{k,t}^{\mathrm{idx}}
      +B_{k,t}^{\mathrm{hdr}}}{B_{k,t}^{\mathrm{raw}}}.
  \end{aligned}
\end{equation}
式~\eqref{eq:effective-compression}第一行中，\(d\) 是原始元素总数，零范数 \(\lVert\cdot\rVert_0\) 统计保留的非零元素数，\(\sigma_{k,t}\) 表示被省略元素的\mbox{比例}。第二行比较实际编码量与未压缩载荷：\(B^{\mathrm{val}}\) 是保留数值的比特，\(B^{\mathrm{idx}}\) 是位置索引，\(B^{\mathrm{hdr}}\) 是码长、形状等头部信息，\(B^{\mathrm{raw}}\) 是原始载荷。有效传输比 \(\rho_{k,t}^{\mathrm{eff}}\) 越小，真实节省越多。稀疏化必须发送索引，否则服务器无法知道每个数值对应哪个参数；因此理论稀疏率高，不等于真实传输比一定同比降低。

误差反馈用于补偿暂未发送的信息。客户端把上一轮残差加入当前更新，再压缩上传；原更新与发送近似值之差留在本地，下一轮继续参与计算。被舍弃的小更新因而能在后续补偿。客户端为此需要保存残差并增加筛选或编码计算；终端频繁退出时，残差状态也难以连续维护\cite{Sattler2020STC,Han2020AdaptiveSparsification,Tian2022GradientCompression}。

代表性工作体现了从单一压缩到自适应压缩的演进。FedPAQ把周期平均、部分参与和量化放入同一框架；一比特压缩感知继续降低数值位宽，但需要稀疏先验和服务器重构\cite{Reisizadeh2020FedPAQ,Zhang2022OneBitCS}。稀疏三值压缩结合双向稀疏、低比特表示和误差累积，自适应方法则按训练状态调整保留比例；近期工作还同时选择稀疏度、量化级和能耗目标\cite{Sattler2020STC,Han2020AdaptiveSparsification,Tian2022GradientCompression,Marnissi2024AdaptiveSQ,Chen2025EdgeAssistedSparsification}。这些方法直接减少式~\eqref{eq:total-bits}中的载荷，而不是加快既定比特的传输。

压缩适合高维模型、窄上行且终端能够承担编码的场景。对小模型、短消息或弱算力设备，索引、头部和编解码成本可能抵消收益。压缩过强会增加优化误差，压缩感知或安全聚合还会引入重构、码本及额外交互\cite{Chen2022SecureCompressedFL}。评价时应同时报告累计比特、精度或损失、编解码成本，并说明是否计入索引。压缩回答“每个客户端发送多少”，空中计算则改变“多个客户端如何同时发送并聚合”。
